\section{RACE-006: Plugin Install Not Locked}
\label{sec:race-006}
\subsection*{Metadata}
\begin{tabular}{ll}
\textbf{Issue ID} & RACE-006 \\
\textbf{Severity} & MEDIUM \\
\textbf{Category} & Race Condition \\
\textbf{Branch} & \branchref{fix/RACE-006-plugin-install-race}
\textbf{Changes} & \branchdiff{fix/RACE-006-plugin-install-race} \\
\end{tabular}

\subsection*{Executive Summary}
The \code{installPlugin} function in \srcfile{src/plugins.ts} performs git clone, directory creation, and file existence checks without any locking. If two concurrent \code{gitclaw plugin install} commands target the same plugin, the race between \code{dirExists} checks, \code{rm} operations, and \code{git clone} calls can result in a corrupted plugin directory — partial clones, overwritten files, or broken \code{node\_modules} state.

\subsection*{Root Cause Analysis}
Two concurrent installs can both pass the directory existence check:

\begin{lstlisting}[language=TypeScript]
// BEFORE — no dedup or lock for concurrent installs
export async function installPlugin(source, targetDir, version?, force?) {
    await mkdir(targetDir, { recursive: true });
    const pluginDir = join(targetDir, name);

    if (await dirExists(pluginDir)) {
        if (await fileExists(join(pluginDir, "plugin.yaml"))) {
            if (force) { await rm(pluginDir, { recursive: true, force: true }); }
            else { return pluginDir; }
        } else { await rm(pluginDir, { recursive: true, force: true }); }
    }

    execFileSync("git", args, { stdio: "pipe" });  // Both processes clone
    return pluginDir;
}
\end{lstlisting}

\subsection*{Fix Implementation}
An in-flight map deduplicates concurrent installs:

\begin{lstlisting}[language=TypeScript]
// AFTER — dedup with in-flight map
const installInFlight = new Map<string, Promise<string>>();

export async function installPlugin(source, targetDir, version?, force?): Promise<string> {
    const key = `${source}::${version || "latest"}`;
    const existing = installInFlight.get(key);
    if (existing && !force) return existing;

    const promise = installPluginInner(source, targetDir, version, force);
    installInFlight.set(key, promise);
    try { return await promise; }
    finally { installInFlight.delete(key); }
}

async function installPluginInner(source, targetDir, version?, force?) {
    // Original implementation
    await mkdir(targetDir, { recursive: true });
    // ... dir checks ...
    execFileSync("git", args, { stdio: "pipe" });
    return pluginDir;
}
\end{lstlisting}

\subsection*{Affected Files}
\begin{itemize}
    \item \srcfile{src/plugins.ts} — installPlugin() split into public + inner function
    \item \srcfile{src/\_\_tests\_\_/plugin-install-race.test.ts}
\end{itemize}

\subsection*{Verification}
Tests verify deduplication of concurrent installs:

\begin{lstlisting}[language=TypeScript]
const [r1, r2] = await Promise.all([
    installPlugin("test-plugin"),
    installPlugin("test-plugin"),
]);
assert.equal(r1, "/plugins/test-plugin");
assert.equal(r2, "/plugins/test-plugin");
assert.equal(callCount.length, 1, "Should only call inner function once");
\end{lstlisting}
